\section{Background}

\subsection{The Self-Healing Concept in Security}

Self-healing systems detect anomalous states and apply remediation
without human intervention. The concept dates to autonomic
computing~\cite{kephart2003vision} and has been operationalised in
commercial products: Microsoft Defender for Endpoint's automated
investigation and response, Tanium's auto-remediation, Tenable's
PVS workflows, IBM SOAR, and Splunk Phantom. The unifying property
is a closed-loop detect--decide--act pipeline whose decisions occur
faster than human review.

\subsection{What Vulnerability Remediation Adds}

Vulnerability remediation is, in principle, an ideal self-healing
target: patches are well-defined, success criteria are observable
(does the host still exhibit the vulnerable configuration?), and
remediation actions are reversible (rollback to pre-patch state).
In practice, self-healing systems in vulnerability management have
been deployed cautiously because of three risks:

\textit{(i) Cascading failures.} A patch that introduces a regression
on one host may affect many hosts if applied widely. Automated
deployment without staged validation can amplify a single faulty
patch into a fleet-wide outage.

\textit{(ii) False prioritisation.} If the scoring rule that selects
patches to apply is biased --- by attacker manipulation
(Paper~4~\S\ref{sec:adversarial}), distributional shift
(Paper~4~\S\ref{sec:external_validity}), or capacity-driven signal
exhaustion (Paper~9~Theorem~1) --- the autonomous system will
remediate the wrong things.

\textit{(iii) Business-criticality gaps.} A self-healing system that
patches a host during business hours, on a customer-facing system,
or without rollback capability causes operational harm
disproportionate to the vulnerability it addresses.

\subsection{Why Now}

Papers 3--9 of this research sequence developed components that
collectively address these risks:

\begin{itemize}
\item HygienePrio (Paper~4)~\cite{paper4hygieneprio} provides a
scoring rule with calibrated weights, real-distribution external
validity (\S\ref{sec:external_validity} of that paper), and
quantified adversarial robustness (\S\ref{sec:adversarial}).
\item Paper 7's lag-1 online calibration~\cite{paper7online} gives
a deployable recipe for moderate capacity ($K \leq 100$).
\item Paper 9's Theorem 1~\cite{paper9selftraj} bounds when
self-driven remediation will exhaust its own signal and identifies
selection-policy decoupling as the structural solution.
\item Real CVE/EPSS/KEV public data
(real\_data/processed/) provides realistic attribute distributions.
\end{itemize}

The motivation for AutoHeal is to integrate these components into
a single framework with explicit safety bounds and to evaluate
honestly whether the result is operationally deployable.

\subsection{Policy Mandates}

CISA Binding Operational Directive 22-01~\cite{cisa2021bod2201}
mandates 15-day remediation of KEV-listed CVEs for federal agencies.
The 2026 DBIR~\cite{verizon2026dbir} reports a 43-day mean time to
remediation across surveyed organisations. A self-healing framework
that materially reduces MTTR has direct policy-compliance value.
NIST SP~800-40 Rev.~4~\cite{nist2022sp80040r4} recommends
risk-based patch management with documented prioritization;
AutoHeal's pre-registered triage rules and frozen evaluation
constitute exactly the form of documentation the standard
contemplates.
